%=====================================================================
\section{Design decisions: what, why, and why not the alternatives}
\label{sec:decisions}

Each decision is stated, then the reason, then the alternatives that were considered and why they were not taken.

\newcommand{\decision}[4]{%
\subsection{#1}#2
\par\noindent\textbf{Decision.} #3
\par\noindent\textbf{Why.} #4}

\decision{D1. BMv2, v1model and a P4Runtime controller}{\label{sec:dec-target}}
{The switch is BMv2's \code{simple\_switch\_grpc} running a P4-16 v1model program; a Python controller talks to it
over P4Runtime.}
{BMv2 is the reference P4 target, runs anywhere, needs no hardware, and supports everything the design uses:
parser loops with \code{lookahead}/\code{advance}, registers, hashes, digests, resubmit, clone, packet-in/out.
P4Runtime is the standard control interface, so ``a high-level program drives the switch'' means exactly: another
P4Runtime client, or this controller's Python API.}
\whynot{\emph{Tofino}: hardware limits (parser depth, stages, no 2,048-bit fields) would dominate the work; the
faculty asked for fingerprinting, not a hardware port. \emph{eBPF/XDP}: not P4. \emph{Other P4 architectures}: v1model is the one
\code{simple\_switch\_grpc} implements completely (resubmit, clone with preserved metadata, digests); PSA support in
BMv2 is a separate, less complete switch, and TNA is Tofino's.}

\decision{D2. The parser reads every fingerprint value}{\label{sec:dec-parser}}
{All variable-length reading (TCP options, TLS lists, HTTP headers, DHCP options, DER) is done by P4 parser states.
The controller never parses a fingerprint field.}
{This is the contribution: fingerprinting in the data plane. The parser is the only P4 block that can loop over a
variable number of items (a state may transition to itself).}
\whynot{\emph{Send every packet to the controller and parse there}: a software fingerprinter with a switch in front.
\emph{Parse with match-action tables}: tables match fixed offsets; variable-length lists need the parser's loops.}

\decision{D3. Results leave the switch as digests}{\label{sec:dec-digest}}
{Each control sends its values in one digest struct (six kinds), as soon as they are complete.}
{A digest carries only the extracted values (a few hundred bytes), asynchronously, without the packet; the packet
is forwarded at once. It is P4Runtime's standard channel for data-plane-to-controller events.}
\whynot{\emph{Packet-in of every interesting packet}: sends the whole packet and makes the controller parse.
\emph{Controller reads registers}: BMv2's P4Runtime server answers UNIMPLEMENTED for register reads (found, see
the problems table), and polling would race with new connections. \emph{Write results into a header and mirror}:
needs a custom receiver and still a clone per message.}

\decision{D4. Sorting, SHA-256 and text in the controller}{\label{sec:dec-hash}}
{The switch sends lists in packet order; the controller sorts the cipher and extension lists, hashes with SHA-256 and
writes FoxIO's text.}
{v1model has no SHA-256 extern, no strings, and no loops in the pipeline (only in the parser). The fingerprint is
fully determined by the values the switch extracted; sorting and hashing them anywhere gives the same result.}
\whynot{\emph{A sorting network in P4} (bitonic, as in the earlier ja4p4 plan): needed on a Tofino to report sorted
lists from the data plane; on BMv2 it adds hundreds of compare-and-swap actions and proves nothing about
fingerprinting. \emph{A CRC instead of SHA-256}: would not be JA4 and could not be compared with FoxIO.}

\decision{D5. Walked packets are resubmitted and forwarded on a second pass}{\label{sec:dec-resubmit}}
{A parser walk skips bytes with \code{advance}; the packet is then resubmitted, and on the second pass only
Ethernet is parsed, so the original bytes are forwarded.}
{\code{advance} removes the skipped bytes from what the deparser emits. Resubmit runs the original packet through
the pipeline again, with chosen metadata preserved (field list 0: arrival port, the pending handshake record).}
\whynot{\emph{Extract every byte into headers}: thousands of header-stack entries; the deparser would re-emit them
(the TCP option walk does extract, since options are at most 40 bytes). \emph{Recirculate}: passes through egress
and the deparser, which would emit the shortened packet.}

\decision{D6. Split and out-of-order messages are joined by the controller and read by the switch}{\label{sec:dec-join}}
{The switch decides which segments may hold pieces of a message (Fig.~\ref{fig:copyrules}) and copies them with
their stream offset; the controller orders the bytes by framing only and sends the complete message back; the
same parser reads it.}
{A P4 switch processes one packet at a time and cannot hold a packet back or keep kilobytes of payload per
connection. Joining needs a buffer; reading does not. So the buffer is in the controller and the reading stays in
P4. The controller's join is framing only (record headers, the empty line, DER lengths), never a fingerprint field.}
\whynot{\emph{Buffer segments in registers}: a ClientHello reaches 2 KB, certificates many KB; registers sized for
that per connection are impossible on hardware and costly on BMv2, and register access is per packet, fixed-width.
\emph{Parse each piece separately and merge partial states}: a list value can be cut in the middle and the order of
pieces is unknown; the state machine would need to be resumable at every byte in both orders. \emph{Ignore split
messages}: wrong by omission for exactly the cases the faculty named.}

\decision{D7. Copies are made at egress (E2E clone)}{\label{sec:dec-e2e}}
{\code{clone\_preserving\_field\_list(CloneType.E2E, 100, 1)} in egress; the clone gets the \code{packet\_in}
header.}
{An ingress-to-egress clone is parsed again from the start, and its walk consumes the bytes again: the copy
arrived without payload (found by the tests). At egress the packet has all its bytes.}

\decision{D8. Per-connection state in hashed register slots}{\label{sec:dec-slots}}
{ConnKey orders the endpoints, hashes them (CRC16 slot, CRC32 check); every control keeps its state in register
arrays of 65,536 entries indexed by the slot, together with the connection's number (D21).}
{A P4 program cannot add table entries; registers are the data plane's writable memory. Ordering the endpoints
makes both directions find the same slot. The 32-bit check tells two connections in one slot apart.}
\whynot{\emph{One table entry per connection from the controller}: a round trip before the first packet's state
exists, and the SYN's timestamp would be lost. \emph{Bigger slots / buckets}: unneeded for one link; one constant.}

\decision{D9. One \code{digest()} call per control; the handshake digest sent last}{\label{sec:dec-onedigest}}
{Every control has exactly one digest call; the handshake record is stored in metadata and sent from a single call
at the end of ingress, on the second pass if this pass sent a message digest.}
{p4c makes one learn list per call site (a second call site of the same struct is a different list the controller
did not enable), and BMv2 keeps \emph{one digest per packet}: the handshake's ACK can carry a whole ClientHello, and
its JA4L would be lost (found with SSH, then covered by \code{test\_one\_digest\_per\_packet}).}

\decision{D10. A serial number in DHCP, HTTP and SSH digests}{\label{sec:dec-serial}\label{sec:serial}}
{A 32-bit counter register per digest kind; each digest carries its value.}
{BMv2 keeps every digest until it is acknowledged and drops a new one identical to it. Identical messages close
together (DHCP retransmissions and their answers, the same HTTP request repeated, identical SSH windows of a bulk
transfer) were lost. 32 bits: the serial need only differ among digests unacknowledged at the same time (rate
$\times$ acknowledgement time, below $2^{20}$ even at $10^6$/s for 1\,s); wrapping at $2^{32}$ is harmless and needs
no flag, because the controller never reads the value.}
\whynot{\emph{Acknowledge faster}: the filter exists until the acknowledgement arrives, however fast. \emph{Use the
timestamp}: two packets in one microsecond collide. \emph{Per-connection counters}: DHCP has no connection.}

\decision{D11. GREASE is recognised in the parser's \code{select}}{}
{\code{IS\_GREASE(v)}: low nibbles of both bytes are 0xA and the two bytes are equal; the value is skipped before
it is stored.}
{GREASE is a bit pattern, not a lookup; a \code{select} on it costs nothing and keeps the lists clean, so counts
are right without a second pass.}
\whynot{\emph{A match-action table marking slots invalid}: needs a list of slots to filter afterwards and then
compaction; the parser simply never stores them.}

\decision{D12. Capacities: 128 / 64 / 64 values, 512 / 1,024 / 256 bytes, then ``incomplete''}{}
{Each list is kept in fixed-width fields (at most 2,048 bits each, chained); past its capacity the count stays
right, an overflow bit is set, and the fingerprint is reported as incomplete.}
{A wrong fingerprint is worse than none. The capacities were checked by a measurement study (every one of 1,238
measured lists fits; closest: cookies 863 of 1,024 bytes, ciphers 86 of 128) and the overflow path by
\code{test\_capacity\_limits}.}
\whynot{\emph{Larger capacities}: cost digest bits and parser iterations for no measured case. \emph{Truncate
silently}: would give wrong hashes.}

\decision{D13. JA4L timed in the switch}{}
{The switch stores the SYN and SYN-ACK times and TTLs and computes both halves when the ACK passes.}
{The point of JA4L is the network position of the measuring device; the switch is that device. The controller
would see digests later and with its own jitter.}
\whynot{\emph{Controller time}: measures the control channel, not the network. Limit kept honest: BMv2 stamps when
its pipeline takes the packet up, so the latency is reported as a measurement and the arithmetic is tested with
known gaps.}

\decision{D14. SSH: counts in the switch, the mode in the controller}{}
{The switch counts SSH packets and bare ACKs per direction and window exactly, and sends the payload lengths in
batches of 8; the controller computes the most frequent length.}
{A mode needs a histogram over arbitrary lengths; a data plane has no dynamic map. Everything else of JA4SSH is
counting, which registers do exactly.}
\whynot{\emph{A fixed histogram in registers}: lengths range up to 64 KB; a per-connection histogram is
impossible, and approximate buckets would give a different fingerprint.}

\decision{D15. The last SSH window through a flush packet}{}
{At the end, the controller sends a packet-out of kind 2 with the connection's endpoints; the switch reports the
partial window itself.}
{FoxIO's tool reports the partial last window; BMv2's P4Runtime cannot read registers.}

\decision{D16. JA4X: the controller cuts certificates out, the switch walks DER}{}
{The controller follows TLS framing to the Certificate message and sends each DER certificate back (kind 3, with its
index); \code{parser\_x509.p4} reads the OIDs.}
{Certificates always span segments, so they need the join path anyway; the DER walk (the fingerprint's reading) is
in P4, with one shared TLV reader and return codes because a parser has no calls.}

\decision{D17. Byte-per-state HTTP walk with chained 2,048-bit fields}{}
{The request is read one byte per parser state; lists grow through chains of 2,048-bit fields by constant-width
slices.}
{HTTP has no length fields: the end of each name, value and line is found byte by byte. BMv2 allows at most 2,048
bits per field and only short variable shifts, so a list longer than 256 bytes must be a chain grown by constant
slices.}

\decision{D18. A rootless replay bench next to the live network}{}
{\code{tools/replay.py} runs the real switch and controller with BMv2's nanomsg packet socket as its ports.}
{Hundreds of runs (every capture, every split case) without root or Mininet; the same binary and controller as in
the live run. The live Mininet run then confirms the result on real TCP stacks.}

\decision{D19. FoxIO's own tools as the reference}{}
{The Rust \code{ja4} tool for most fingerprints, tshark fields for JA4TS (not in the open tool), FoxIO's Wireshark
plugin output (or the same rule on tshark fields) for JA4D.}
{Independent of this project: a match with the authors' implementation is the strongest evidence of correctness.}
\whynot{\emph{Compare with my own Python implementation}: two implementations of my own reading could share the
same misunderstanding.}

\decision{D20. One design only}{}
{The two approaches discussed at the start (pure P4 without external programs; P4 invoking high-level code) were
set aside by the faculty; the design is one switch program plus one controller.}
{Focus on correct fingerprints first.}

\decision{D21. A number for every connection (port reuse)}{\label{sec:dec-connid}}
{Control \code{ConnId} (\code{main.p4}) gives each connection a number from a counter: when another connection
comes to the slot (a different check), and when a SYN with a new initial sequence number arrives on the same
endpoints. Every control keeps its state with this number instead of the check; every TCP digest and every copy
carries it; the controller keys records, the reassembler and SSH windows by (endpoints, number); the comparison keys
FoxIO's records by (endpoints, place among the TCP streams on those endpoints).}
{Endpoints do not identify a connection: clients and NATs reuse ports. Before, the controller kept the first value
per endpoints (the second connection's JA4 was lost) and the SSH counts of two connections were added together into
a wrong value. A repeated SYN (same sequence number) must not start a new connection.}
\whynot{\emph{Key by the client's initial sequence number in the controller}: only the SYN carries it; every
digest would need it, and connections without a SYN seen have none. \emph{Detect a new connection in the controller
from a new JA4T}: digests of different kinds arrive in no fixed order. \emph{Key on FIN / RST}: they can be lost
or missing from a capture; a SYN with a new sequence number is what tshark and FoxIO use.}

\decision{D22. Rule (e): pieces ahead of a gap in an HTTP direction}{\label{sec:dec-gap}}
{Per direction the switch keeps the edge (after the furthest byte seen). In a direction that has carried an HTTP
request, a segment beyond the edge opens a gap, and segments within 16 KB of the gap's start are copied until a
segment covering the gap's first byte arrives. The controller lists a connection's JA4H by each request's sequence
number.}
{Rule (c) covers only a direction's first message. On a keep-alive connection the second request, cut in two and
second piece first, was never reported. Reporting in stream order is FoxIO's order; a joined request otherwise
arrives after the whole requests that followed it.}
\whynot{\emph{Copy every segment after a gap in every connection}: TLS and SSH connections carry nothing more to
read after their hellos, so copies there are only load. \emph{Track every hole}: one gap per direction covers the
case at the cost of four registers; further holes inside a pending message are covered by rule (b).}

\decision{D23. VLAN tags and SSH on any port}{}
{The parser reads up to two VLAN tags (802.1Q, 802.1ad outer). A connection is SSH when it is on port 22 or a side's
data begins with ``SSH-''; bare ACKs are counted from the start of every connection, so those before the banner
count.}
{FoxIO's tool reads through VLAN tags and recognises SSH by its banner on any port (tshark's heuristic): the switch
reported nothing on tagged frames and on SSH on port 2222. FoxIO counts the handshake's ACK of such a connection as
a client bare ACK (the first version gave \code{c6s6} for FoxIO's \code{c7s6}).}
\whynot{\emph{A list of SSH ports from the controller}: misses the unknown ones, and FoxIO needs no list.
\emph{More than two tags}: not used on access links; such frames are forwarded, not fingerprinted.}

\decision{D24. Bounded controller buffers}{}
{The reassembler keeps at most 32 KB per direction (lowest offsets dropped first), drops a direction idle for 60 s,
and keeps at most 100,000 directions.}
{Without bounds a long run kept every copied piece forever, and half-sent messages could fill the controller.}
\whynot{\emph{Drop a message's bytes as soon as it is joined}: tried, and it broke 3 of the 202 split cases. A
one-byte segment in the middle of a request that begins like a request is ``joined'' first; the real request, whose
first segment comes last, then lacked those bytes.}
